\documentclass{article}

% Evaluations & Datasets Track, double-blind
\usepackage[eandd]{neurips_2026}

\usepackage[utf8]{inputenc}
\usepackage[T1]{fontenc}
\usepackage{hyperref}
\usepackage{url}
\usepackage{booktabs}
\usepackage{amsfonts}
\usepackage{nicefrac}
\usepackage{microtype}
\usepackage{xcolor}
\usepackage{enumitem}
\usepackage{multirow}
\usepackage{graphicx}

\title{AgentTelemetry: A Grounded Taxonomy and Empirical Evaluation\\of Observability Primitives for Autonomous AI Agent Systems}

\author{Anonymous Authors}

\begin{document}

\maketitle

\begin{abstract}
Autonomous AI agents built on large language models are transitioning
from research prototypes to production deployments, yet the
observability infrastructure needed to diagnose their failures remains
fundamentally inadequate.  OpenTelemetry's GenAI semantic conventions
cover LLM invocation, tool execution, and retrieval, but leave five
critical agent orchestration
phases---planning, reasoning, safety-check monitoring, inter-agent delegation,
and memory management---without any span-level representation.
We present \textsc{AgentTelemetry}, an open-source, OTel-native
observability library that introduces a \emph{grounded taxonomy} of
nine agent-specific span kinds derived from systematic analysis of
seven production agent frameworks.
We make four contributions: (1)~the taxonomy itself, validated through
an ablation study showing every span kind is necessary for detecting
at least one fault type; (2)~a case study on 112~SWE-bench Lite
instances revealing that reasoning loops account for 75\% of agent
failures (95\% CI: [66\%, 82\%])---a failure mode invisible to vanilla
OTel---and a telemetry-guided intervention improving the patch rate by
$+11$~pp (3-seed mean, $p{<}0.05$); (3)~structural completeness
via controlled fault injection across 14~fault types, 7~frameworks,
and 5~observability conditions (FDR\,=\,1.000 vs.\ 0.429 for vanilla
OTel); and (4)~a pip-installable library with adapters for seven
frameworks.  Instrumentation overhead is 11.7\,\textmu s median per
span, negligible relative to LLM API latencies.
\end{abstract}

%% ====================================================================
%% 1. INTRODUCTION
%% ====================================================================
\section{Introduction}
\label{sec:introduction}

Large language model (LLM) based autonomous agents represent a paradigm
shift in software architecture.  Unlike conventional services that
execute deterministic logic, an LLM agent pursues a goal through
iterative reasoning, planning, tool use, and
self-reflection---frequently delegating sub-tasks to other
agents~\citep{agentbench,agentboard}.  Frameworks such as LangChain,
CrewAI, AutoGen, and LlamaIndex have accelerated adoption, yet
production teams cannot reliably answer fundamental diagnostic
questions: \emph{Why did the agent choose this tool? Where did the
hallucination originate? How much did this task cost across all agents?}

OpenTelemetry (OTel)~\citep{otel} has become the standard for
distributed tracing.  Its GenAI semantic conventions define operation
types for LLM calls, retrieval, tool execution, and agent lifecycle.
However, five orchestration phases---planning, reasoning, safety-check
monitoring, inter-agent delegation, and memory management---have no
representation in OTel GenAI, leaving them as opaque
\texttt{INTERNAL} spans indistinguishable from any other application
logic.

We present \textsc{AgentTelemetry}, an open-source, OTel-native
library that fills this gap with four contributions:

\begin{enumerate}[leftmargin=*,nosep]
\item \textbf{A grounded taxonomy of nine agent-specific span kinds},
  derived from seven production frameworks and validated through
  ablation (\S\ref{sec:taxonomy}).

\item \textbf{A SWE-bench case study demonstrating actionable
  observability:} reasoning loops account for 75\% of failures on
  112~instances, and a telemetry-guided intervention improves the patch
  rate by $+11$~pp (\S\ref{sec:rq6}).

\item \textbf{Structural completeness via controlled fault injection:}
  FDR\,=\,1.000 across 14~fault types versus 0.429 for vanilla
  OTel (\S\ref{sec:evaluation}).

\item \textbf{An open-source, pip-installable library} (3,700+ LOC,
  78~tests) with seven framework adapters (\S\ref{sec:implementation}).
\end{enumerate}

%% ====================================================================
%% 2. BACKGROUND
%% ====================================================================
\section{Background and Motivation}
\label{sec:background}

An LLM agent's execution produces a deeply nested, non-deterministic
trace of cognitive operations.  A multi-agent research
assistant---comprising planner, researcher, and writer agents---may
produce dozens of spans, accumulating thousands of tokens and
non-trivial cost.  Debugging a failure requires tracing backward
across agent boundaries; without unified tracing, operators resort to
ad-hoc logging that is inconsistent across frameworks and scales
poorly.

The OTel GenAI semantic conventions define eight
\texttt{gen\_ai.operation.name} values---\texttt{chat},
\texttt{text\_completion}, \texttt{embeddings}, \texttt{retrieval},
\texttt{execute\_tool}, \texttt{generate\_content},
\texttt{create\_agent}, and \texttt{invoke\_agent}---covering LLM
invocation, vector retrieval, tool execution, and agent lifecycle.
Despite this breadth, five of nine \textsc{AgentTelemetry} span
kinds---\texttt{PLANNING}, \texttt{REASONING}, \texttt{GUARD\_RAIL},
\texttt{DELEGATION}, and \texttt{MEMORY}---have no OTel GenAI
equivalent (Table~\ref{tab:gap}).

\begin{table}[t]
\centering
\small
\caption{Agent concepts mapped to observability primitives.
$\bullet$\,=\,overlapping with OTel GenAI;
$\triangleright$\,=\,extends OTel GenAI; $\star$\,=\,novel.}
\label{tab:gap}
\begin{tabular}{@{}llll@{}}
\toprule
\textbf{Agent Concept} & \textbf{OTel GenAI} & \textbf{AgentTelemetry} & \\
\midrule
Agent lifecycle    & create/invoke\_agent  & \textbf{AGENT} & $\triangleright$ \\
LLM invocation     & chat / text\_compl.   & \textbf{LLM\_CALL} & $\bullet$ \\
Tool execution     & execute\_tool         & \textbf{TOOL\_CALL} & $\bullet$ \\
Doc.\ retrieval    & retrieval, embeddings & \textbf{RETRIEVAL} & $\bullet$ \\
Task decomposition & (none)                & \textbf{PLANNING} & $\star$ \\
Chain-of-thought   & (none)                & \textbf{REASONING} & $\star$ \\
Safety monitoring  & (none)                & \textbf{GUARD\_RAIL} & $\star$ \\
Inter-agent handoff & (none)               & \textbf{DELEGATION} & $\star$ \\
Memory read/write  & (none)                & \textbf{MEMORY} & $\star$ \\
\bottomrule
\end{tabular}
\end{table}

%% ====================================================================
%% 3. SPAN TAXONOMY
%% ====================================================================
\section{A Grounded Taxonomy of Agent Span Kinds}
\label{sec:taxonomy}

\subsection{Derivation Methodology}

The nine span kinds were derived through a structured qualitative
analysis modeled on open coding~\citep{strauss_corbin}.  The process
proceeded in four stages.

\textbf{Stage~1: API inventory.}  For each of seven frameworks
(LangChain, CrewAI, AutoGen, LlamaIndex, Anthropic~SDK, OpenAI~SDK,
Custom), we catalogued every public API entry point representing a
distinct execution phase.  Table~\ref{tab:traceability} lists the
25~entry points examined.

\textbf{Stage~2: Open coding.}  Each entry point was assigned a
candidate label describing its execution phase, producing 14~initial
candidates.

\textbf{Stage~3: Axial coding and merging.}  We merged candidates
using two criteria: (a)~\emph{observational
indistinguishability}---candidates that cannot be reliably
distinguished from framework-emitted signals are merged; and
(b)~\emph{diagnostic redundancy}---candidates triggering the same
fault-detection logic are merged.  Five merges occurred:
\textsc{LLM\_Chat}+\textsc{LLM\_Completion}$\to$\texttt{LLM\_CALL};
\textsc{Reflection}$\to$\texttt{REASONING};
\textsc{Agent\_Comm}$\to$\texttt{DELEGATION};
\textsc{Read/Write\_Memory}$\to$\texttt{MEMORY} with
\texttt{memory.operation} attribute;
\textsc{Human\_Input}$\to$removed (modeled as
\texttt{AGENT} with \texttt{agent.type=``human''}).

\textbf{Stage~4: Saturation check.}  After five frameworks, the set
had stabilized at nine.  Two more frameworks introduced no new
candidates.  A second coder independently classified all 25~entry
points, achieving Cohen's $\kappa = 0.904$ (almost perfect agreement).

\begin{table}[t]
\centering
\scriptsize
\caption{Traceability matrix: 25~API entry points across 7~frameworks
mapped to 9~span kinds (subset shown for space).}
\label{tab:traceability}
\begin{tabular}{@{}llll@{}}
\toprule
\textbf{Framework} & \textbf{API Entry Point} & \textbf{Candidate} & \textbf{Final Kind} \\
\midrule
\multirow{4}{*}{LangChain}
  & \texttt{on\_chain\_start} (agent) & Agent & AGENT \\
  & \texttt{on\_llm\_start} & LLM\_Chat & LLM\_CALL \\
  & \texttt{on\_tool\_start} & Tool\_Call & TOOL\_CALL \\
  & \texttt{on\_retriever\_start} & Retrieval & RETRIEVAL \\
\midrule
AutoGen & \texttt{BaseChatAgent.run()} & Agent & AGENT \\
\midrule
LlamaIndex & \texttt{span\_enter} (sub\_question) & Planning & PLANNING \\
\midrule
\multirow{3}{*}{Custom}
  & \texttt{start\_guardrail\_span()} & Guard\_Rail & GUARD\_RAIL \\
  & \texttt{start\_delegation\_span()} & Delegation & DELEGATION \\
  & \texttt{start\_memory\_span()} & Memory & MEMORY \\
\bottomrule
\end{tabular}
\end{table}

\subsection{The Nine Span Kinds}

All nine kinds are represented as semantic attributes on standard OTel
\texttt{INTERNAL} spans via \texttt{agenttelemetry.span.kind},
ensuring OTLP export works with any compatible backend.  Of the nine
(Table~\ref{tab:spantaxonomy}), three overlap with OTel GenAI, one
extends it, and \textbf{five are novel}: \texttt{PLANNING},
\texttt{REASONING}, \texttt{GUARD\_RAIL}, \texttt{DELEGATION}, and
\texttt{MEMORY}.

\begin{table}[t]
\centering
\small
\caption{The nine agent-specific span kinds with key attributes.
$\star$\,=\,novel (no OTel GenAI equivalent).}
\label{tab:spantaxonomy}
\begin{tabular}{@{}llp{3.8cm}@{}}
\toprule
\textbf{\#} & \textbf{Span Kind} & \textbf{Key Attributes} \\
\midrule
1 & \texttt{AGENT}      & \texttt{agent.name}, \texttt{agent.framework}, \texttt{agent.role} \\
2 & \texttt{LLM\_CALL}  & \texttt{llm.model}, \texttt{llm.input\_tokens}, \texttt{llm.output\_tokens}, \texttt{llm.cost} \\
3 & \texttt{TOOL\_CALL} & \texttt{tool.name}, \texttt{tool.input}, \texttt{tool.output}, \texttt{tool.status} \\
4 & \texttt{PLANNING}\,$\star$   & \texttt{planning.strategy}, \texttt{planning.step\_count} \\
5 & \texttt{REASONING}\,$\star$  & \texttt{reasoning.chain} \\
6 & \texttt{RETRIEVAL}  & \texttt{retrieval.source}, \texttt{retrieval.query}, \texttt{retrieval.doc\_count} \\
7 & \texttt{GUARD\_RAIL}\,$\star$& \texttt{guardrail.name}, \texttt{guardrail.result} \\
8 & \texttt{DELEGATION}\,$\star$ & \texttt{delegation.source\_agent}, \texttt{delegation.target\_agent} \\
9 & \texttt{MEMORY}\,$\star$     & \texttt{memory.operation}, \texttt{memory.key} \\
\bottomrule
\end{tabular}
\end{table}

\textbf{PLANNING vs.\ REASONING.}
\texttt{PLANNING} spans wrap LLM calls producing \emph{structured
action plans} (task decompositions, sub-query lists).
\texttt{REASONING} spans wrap \emph{chain-of-thought steps within a
single action} (intermediate inference, reflection).  The ablation
(\S\ref{sec:rq2}) confirms both are needed: removing
\texttt{PLANNING} makes planning-failure faults undetectable; removing
\texttt{REASONING} makes reasoning-loop faults undetectable.

\textbf{Why nine and not fewer?}  Merging \texttt{PLANNING} into
\texttt{REASONING} would lose the ability to distinguish task
decomposition (``break this into 3~sub-queries'') from chain-of-thought
(``the answer is X because Y'').  In ReAct-style agents, the adapter
creates one \texttt{LLM\_CALL} span for the combined output; the parent
is \texttt{PLANNING} when the output contains an \texttt{Action:} block,
and \texttt{REASONING} when it contains only \texttt{Thought:} blocks.

\textbf{Why nine and not more?}  We considered splitting
\texttt{LLM\_CALL} into chat and completion variants, but the same
detection logic (token overflow, cost explosion, hallucination) applies
to both.

\subsection{Cross-Framework Validation}

No single framework natively exposes all nine execution phases.
Table~\ref{tab:framework_coverage} shows the mapping.  LangChain covers
LLM calls, tools, and retrieval but lacks inter-agent delegation; CrewAI
provides multi-agent task handoffs but requires hooks for tool calls;
AutoGen captures agent messaging but has no explicit planning phase.
The taxonomy is framework-agnostic by construction: it captures
\emph{what} the agent does, not \emph{how}.

\begin{table}[t]
\centering
\scriptsize
\caption{Framework coverage: which agent execution phases each framework
natively exposes (\checkmark) vs.\ requires adapter mapping ($\circ$).}
\label{tab:framework_coverage}
\begin{tabular}{@{}lccccccc@{}}
\toprule
\textbf{Span Kind} & \rotatebox{70}{\textbf{LangChain}} &
\rotatebox{70}{\textbf{CrewAI}} & \rotatebox{70}{\textbf{AutoGen}} &
\rotatebox{70}{\textbf{LlamaIndex}} & \rotatebox{70}{\textbf{Anthropic}} &
\rotatebox{70}{\textbf{OpenAI}} & \rotatebox{70}{\textbf{Custom}} \\
\midrule
AGENT       & \checkmark & \checkmark & \checkmark & \checkmark & $\circ$ & $\circ$ & \checkmark \\
LLM\_CALL   & \checkmark & \checkmark & \checkmark & \checkmark & \checkmark & \checkmark & \checkmark \\
TOOL\_CALL  & \checkmark & $\circ$    & \checkmark & \checkmark & \checkmark & \checkmark & \checkmark \\
PLANNING    & $\circ$    & $\circ$    & $\circ$    & \checkmark & $\circ$ & $\circ$ & \checkmark \\
REASONING   & \checkmark & \checkmark & \checkmark & $\circ$    & $\circ$ & $\circ$ & \checkmark \\
RETRIEVAL   & \checkmark & $\circ$    & $\circ$    & \checkmark & $\circ$ & $\circ$ & \checkmark \\
GUARD\_RAIL & $\circ$    & $\circ$    & $\circ$    & $\circ$    & $\circ$ & $\circ$ & \checkmark \\
DELEGATION  & $\circ$    & \checkmark & \checkmark & $\circ$    & $\circ$ & $\circ$ & \checkmark \\
MEMORY      & $\circ$    & $\circ$    & $\circ$    & $\circ$    & $\circ$ & $\circ$ & \checkmark \\
\bottomrule
\end{tabular}
\end{table}

\textbf{Cross-validation with competing taxonomies.}
AgentDebug's five competency categories~\citep{agentdebug} map directly:
memory$\to$\texttt{MEMORY}, reflection$\to$\texttt{REASONING},
planning$\to$\texttt{PLANNING}, action$\to$\texttt{TOOL\_CALL},
system-level$\to$\texttt{AGENT}.
MAST's 14~failure modes~\citep{mast} operate at a different
abstraction level, but 12 of~14 (85.7\%) are detectable via our span
kinds (Table~\ref{tab:mast-mapping} in Appendix).  The two
gaps---failure to ask for clarification (FM-2.2) and missing
verification (FM-3.2)---are \emph{omission} failures requiring
expected-span specifications beyond trace-level anomaly detection.

%% ====================================================================
%% 4. IMPLEMENTATION
%% ====================================================================
\section{Implementation}
\label{sec:implementation}

\textsc{AgentTelemetry} is a pip-installable Python package (3,700+
LOC, 78~tests, Apache-2.0) built on the OpenTelemetry SDK.  The
library comprises three modules: \textbf{Core} (span definitions,
privacy filtering, context propagation, exporters),
\textbf{Adapters} (seven framework-specific instrumentors), and
\textbf{Analysis} (decision attribution, hallucination tracing, cost
aggregation, anomaly detection).

Three adapter strategies cover all seven frameworks:
\emph{callback handlers} (LangChain, CrewAI, LlamaIndex; 903~LOC),
\emph{monkey-patching} (Anthropic SDK, OpenAI SDK, AutoGen; 795~LOC),
and \emph{manual API} (Custom; 315~LOC).  All framework imports occur
inside \texttt{\_instrument()}, avoiding import-time dependencies.

A \texttt{PrivacyLevel} enum controls attribute filtering:
\texttt{NONE} (structural only), \texttt{METADATA\_ONLY} (+ model
name, token counts, costs), and \texttt{FULL} (+ prompt/completion
text).  The default is \texttt{METADATA\_ONLY}.

\textbf{Telemetry-driven runtime control.}
The \texttt{AgentCircuitBreaker}, an OTel \texttt{SpanProcessor},
intercepts spans as they complete and checks for fault patterns in
real time.  Four policies are supported: reasoning loop ($\geq N$
identical tool calls), cost explosion, context overflow, and
delegation cycle.  This mechanism is impossible without agent-specific
span kinds: the circuit breaker dispatches on
\texttt{agenttelemetry.span.kind}.

%% ====================================================================
%% 5. EVALUATION
%% ====================================================================
\section{Evaluation}
\label{sec:evaluation}

We evaluate through six research questions:

\noindent
\begin{tabular}{@{}lp{5.8cm}@{}}
\textbf{RQ1} & Does AgentTelemetry improve fault detection over vanilla OTel? \\
\textbf{RQ2} & Are all nine span kinds necessary? \\
\textbf{RQ3} & What is the instrumentation overhead? \\
\textbf{RQ5} & Does AgentTelemetry work with real LLM APIs? \\
\textbf{RQ6} & Does AgentTelemetry enable diagnosis and remediation of real failures? \\
\end{tabular}

\subsection{Experimental Setup}

\textbf{Task.}  A standardized research assistant agent that receives a
question, plans search queries, retrieves documents, calls tools,
reasons over results, monitors safety-check outcomes, and generates an
answer.  This exercises all nine span kinds.

\textbf{Models.}  Six LLMs across three capability tiers and two
providers: Haiku~4.5, Sonnet~4.6, Opus~4.6 (Anthropic); GPT-4o-mini,
O3-mini, GPT-4o (OpenAI).

\textbf{Frameworks.}  Seven adapters: LangChain, CrewAI, AutoGen,
LlamaIndex, Anthropic~SDK, OpenAI~SDK, and Custom.  The Custom adapter
serves as the \emph{reference implementation} with full fault coverage
across all 14~fault types.  The remaining six are adapter-validated
stubs that exercise the instrumentation layer through manual
instrumentation, validating adapter correctness without requiring
framework installation.

\textbf{Reproducibility.}  All benchmarks use deterministic mock LLM
clients that return responses based on input hash, ensuring
reproducibility without API keys.  The fault injector records ground
truth for every injected fault.

\subsection{RQ1: Fault Detection Rate}
\label{sec:rq1}

The controlled benchmark measures \emph{structural completeness}:
does the span taxonomy provide a sufficient detection signal for
every fault type?  FDR\,=\,1.000 is an \emph{upper bound}, not a
field detection rate.

\textbf{Method.}  14~fault types spanning all nine span kinds
(Table~\ref{tab:faults}) are injected at rate~1.0.  Five observability
conditions: no telemetry, vanilla OTel, OTel+GenAI,
\textsc{AgentTelemetry} metadata-only, and full capture.
14$\times$5$\times$7$\times$6 = 2,940 configurations, zero errors.

\begin{table}[t]
\centering
\small
\caption{14~fault types and the span kind required for detection.}
\label{tab:faults}
\begin{tabular}{@{}rlll@{}}
\toprule
\textbf{\#} & \textbf{Fault Type} & \textbf{Detection Signal} & \textbf{Required Kind} \\
\midrule
1  & Wrong tool          & Ground truth mismatch  & TOOL\_CALL \\
2  & Hallucination       & No retrieval grounding & LLM\_CALL \\
3  & Infinite loop       & $\geq$3 identical calls & TOOL\_CALL \\
4  & Context overflow    & Token growth $>$1.3$\times$ & LLM\_CALL \\
5  & Cost explosion      & Cost $>$\$0.10        & LLM\_CALL \\
6  & Circular delegation & A$\to$B$\to$A cycle    & DELEGATION \\
7  & Tool failure        & Error status           & TOOL\_CALL \\
8  & Timeout             & Error + timeout msg    & LLM\_CALL \\
9  & Stale retrieval     & Staleness $>$3600s     & RETRIEVAL \\
10 & Guardrail bypass    & Result = ``bypass''    & GUARD\_RAIL \\
11 & Planning failure    & Steps $>$10            & PLANNING \\
12 & Reasoning loop      & Identical chains       & REASONING \\
13 & Agent misroute      & Misrouted flag         & AGENT \\
14 & Memory corruption   & Corrupted flag         & MEMORY \\
\bottomrule
\end{tabular}
\end{table}

\textbf{Results.}  Table~\ref{tab:fdr} shows four key findings.
\emph{Finding~1:} Without telemetry, zero faults are detected---agents
fail silently.
\emph{Finding~2:} Vanilla OTel detects 6/14 faults (FDR\,=\,0.429)
using generic span attributes (\texttt{tool.name},
\texttt{llm.input\_tokens}, error status).  The remaining eight faults
require agent-specific span kinds that vanilla OTel does not provide.
\emph{Finding~3:} OTel+GenAI achieves the \emph{same} FDR\,=\,0.429---despite
adding \texttt{gen\_ai.*} attributes, it lacks the five novel span kinds.
The detection gap is structural, not an artifact of detector implementation.
\emph{Finding~4:} \textsc{AgentTelemetry} achieves FDR\,=\,1.000 at both
\texttt{METADATA\_ONLY} and \texttt{FULL} privacy levels, indicating that
metadata-level capture is sufficient for fault detection---full content
capture adds debugging detail but does not change detection rates.

\begin{table}[t]
\centering
\small
\caption{Fault Detection Rate by condition (reference implementation).}
\label{tab:fdr}
\begin{tabular}{@{}lccccc@{}}
\toprule
\textbf{Fault Type} & \textbf{None} & \textbf{V-OTel} & \textbf{GenAI} & \textbf{Meta} & \textbf{Full} \\
\midrule
Wrong tool          & 0.0 & 1.0 & 1.0 & 1.0 & 1.0 \\
Hallucination       & 0.0 & 0.0 & 0.0 & 1.0 & 1.0 \\
Infinite loop       & 0.0 & 1.0 & 1.0 & 1.0 & 1.0 \\
Context overflow    & 0.0 & 1.0 & 1.0 & 1.0 & 1.0 \\
Cost explosion      & 0.0 & 1.0 & 1.0 & 1.0 & 1.0 \\
Circular delegation & 0.0 & 0.0 & 0.0 & 1.0 & 1.0 \\
Tool failure        & 0.0 & 1.0 & 1.0 & 1.0 & 1.0 \\
Timeout             & 0.0 & 1.0 & 1.0 & 1.0 & 1.0 \\
Stale retrieval     & 0.0 & 0.0 & 0.0 & 1.0 & 1.0 \\
Guardrail bypass    & 0.0 & 0.0 & 0.0 & 1.0 & 1.0 \\
Planning failure    & 0.0 & 0.0 & 0.0 & 1.0 & 1.0 \\
Reasoning loop      & 0.0 & 0.0 & 0.0 & 1.0 & 1.0 \\
Agent misroute      & 0.0 & 0.0 & 0.0 & 1.0 & 1.0 \\
Memory corruption   & 0.0 & 0.0 & 0.0 & 1.0 & 1.0 \\
\midrule
\textbf{Aggregate FDR} & \textbf{0.000} & \textbf{0.429} & \textbf{0.429} & \textbf{1.000} & \textbf{1.000} \\
\bottomrule
\end{tabular}
\end{table}

\textbf{Statistical robustness.}  With 5~seeds at 5~injection rates
(10\%--100\%), FDR\,=\,1.000$\pm$0.000 at 100\% injection and
0.921$\pm$0.158 at 10\%.  FPR\,=\,0.000 across all false-positive runs
and separately confirmed on 112~clean SWE-bench traces (3,060~spans).

\subsection{RQ2: Span Kind Necessity (Ablation)}
\label{sec:rq2}

For each span kind, we remove all spans of that kind from the exported
trace and re-run fault detection, producing a 9$\times$14 necessity
matrix.  \emph{Every span kind is required for at least one fault type}
(Table~\ref{tab:ablation}).  The matrix exhibits a clean
block-diagonal structure: LLM\_CALL is most broadly required
(4~faults: hallucination, context overflow, cost explosion, timeout);
TOOL\_CALL covers 2~faults (wrong tool, infinite loop, tool failure);
each of the five novel kinds has exactly one unique fault type that
becomes undetectable without it.

This directly addresses the concern that ``nine span kinds are
introduced without justification.''  The ablation proves necessity: no
kind can be removed without losing detection of at least one failure
mode.  The taxonomy is also extensible---span kinds are string
attributes, so users can define new kinds without modifying the library.

\begin{table}[t]
\centering
\scriptsize
\caption{Ablation matrix: \textbf{R}\,=\,required (FDR drops to 0).
Every span kind is required for $\geq$1 fault.}
\label{tab:ablation}
\begin{tabular}{@{}l*{9}{c}@{}}
\toprule
\textbf{Removed} & \rotatebox{70}{wrong\_tool} & \rotatebox{70}{hallucin.} & \rotatebox{70}{inf.\ loop} & \rotatebox{70}{ctx.\ overflow} & \rotatebox{70}{cost expl.} & \rotatebox{70}{circ.\ deleg.} & \rotatebox{70}{tool fail.} & \rotatebox{70}{timeout} & \rotatebox{70}{\textbf{Unique}} \\
\midrule
AGENT       & -- & -- & -- & -- & -- & -- & -- & -- & misroute \\
LLM\_CALL   & -- & \textbf{R} & -- & \textbf{R} & \textbf{R} & -- & -- & \textbf{R} & \\
TOOL\_CALL  & -- & -- & \textbf{R} & -- & -- & -- & \textbf{R} & -- & \\
PLANNING    & -- & -- & -- & -- & -- & -- & -- & -- & plan.\ fail \\
REASONING   & -- & -- & -- & -- & -- & -- & -- & -- & reas.\ loop \\
RETRIEVAL   & -- & -- & -- & -- & -- & -- & -- & -- & stale ret. \\
GUARD\_RAIL & -- & -- & -- & -- & -- & -- & -- & -- & GR bypass \\
DELEGATION  & -- & -- & -- & -- & -- & \textbf{R} & -- & -- & \\
MEMORY      & -- & -- & -- & -- & -- & -- & -- & -- & mem.\ corr. \\
\bottomrule
\end{tabular}
\end{table}

\subsection{RQ3: Instrumentation Overhead}
\label{sec:rq3}

Across all 9~span kinds ($N{=}90{,}000$ measurements), the median
span creation latency is 11.7\,\textmu s (p95\,=\,13.3\,\textmu s,
p99\,=\,28.2\,\textmu s).  Given LLM API round trips of 500\,ms to
10\,s, instrumentation adds $<$0.006\% overhead.  Memory cost is
$\sim$3.1\,KB per span.  Sustained throughput exceeds
58,000~spans/sec.  Under 100-thread concurrent pressure,
p99\,=\,76.8\,\textmu s; a 1,200-span long-running trace shows only
7.6\% latency degradation.

\subsection{RQ5: Real LLM API Validation}
\label{sec:rq5}

To address the primary threat to validity---that our core evaluation uses
mock LLM clients---we run \textsc{AgentTelemetry} against real API
endpoints from two providers.  We design a multi-hop QA agent with
5~tools and 20~questions across 4~categories.

\textbf{Trace structure.}
Table~\ref{tab:rq5-structure} summarizes results across 8~models
(159~runs).  Each model produces well-formed trace trees with the
expected span hierarchy (AGENT $\to$ PLANNING $\to$ MEMORY $\to$
[REASONING $\to$ LLM\_CALL $\to$ tool spans]$^*$ $\to$ MEMORY).  All
four analysis modules run successfully without modification.

\begin{table}[t]
\centering
\small
\caption{RQ5: Trace structure from real LLM API experiment.
Averages across 20~questions per model.}
\label{tab:rq5-structure}
\begin{tabular}{@{}lrrrrr@{}}
\toprule
\textbf{Model} & \textbf{Runs} & \textbf{Avg It.} & \textbf{Avg Tools} & \textbf{Avg In Tok} & \textbf{Err.} \\
\midrule
GPT-4o-mini     & 20 & 3.2 & 2.9 & 2{,}335 & 0 \\
GPT-4o          & 20 & 3.0 & 2.6 & 2{,}157 & 0 \\
GPT-4.1-nano    & 19 & 3.3 & 3.1 & 2{,}520 & 1 \\
GPT-4.1-mini    & 20 & 3.0 & 2.5 & 2{,}052 & 0 \\
GPT-4.1         & 20 & 3.4 & 2.8 & 2{,}581 & 0 \\
O3-mini         & 20 & 2.6 & 1.6 & 1{,}699 & 0 \\
O4-mini         & 20 & 3.2 & 2.2 & 2{,}168 & 0 \\
Haiku~4.5       & 20 & 3.0 & 2.4 & 4{,}768 & 0 \\
\midrule
\textbf{Total}  & \textbf{159} & \textbf{3.1} & \textbf{2.5} & \textbf{2{,}535} & \textbf{1} \\
\bottomrule
\end{tabular}
\end{table}

\textbf{Natural faults.}
The \texttt{AnomalyDetector} found \emph{organic} faults in clean
traces: GPT-4o-mini and GPT-4.1-nano both triggered
\texttt{infinite\_retry} alerts by calling the same tool 3--4 times
with identical arguments---genuine model behaviors, not injected faults,
demonstrating that \textsc{AgentTelemetry} detects real failure patterns.

\textbf{Overhead.}
Over 15~paired runs, the median wall-clock overhead of instrumentation
is $<$5\%, but individual runs range from $-50\%$ to $+162\%$,
reflecting LLM API latency jitter (1--6\,s round trips), not
instrumentation cost.

\subsection{RQ6: Failure Characterization on SWE-bench}
\label{sec:rq6}

We instrument a coding agent on 112~instances from SWE-bench
Lite~\citep{swebench} (37\% of the benchmark, all 12~repositories).

\textbf{Setup.}  ReAct architecture with 5~tools, 8~iterations max,
GPT-4o-mini.  Total cost: \$2.53.

\textbf{Results.}  The agent produced 3,060~spans.  Of 112~instances,
26 (23\%) produced patches; 84 (75\%) exhausted the iteration limit.

\textbf{Fault-type distribution.}
Two dominant failure modes in the 84~failed instances:
\begin{itemize}[nosep]
\item \textbf{Reasoning loop (75\%, 95\% CI [66\%, 82\%]):} The agent
  repeatedly called \texttt{search\_code} with similar queries without
  converging---visible as $\geq$4 consecutive REASONING$\to$LLM\_CALL
  cycles.  \emph{Invisible to vanilla OTel.}
\item \textbf{Tool failure (15\%):} \texttt{read\_file} returned ERROR
  for files outside the changed set.
\end{itemize}

\textbf{Example trace.}  Instance \texttt{django\_\_django-10914}: the
agent produced the cycle REASONING $\to$ LLM\_CALL $\to$
TOOL\_CALL(\texttt{search\_code}, ``FilePathField'') four consecutive
times, each returning the same limited results.  Without REASONING
spans, these four cycles collapse into eight undifferentiated
\texttt{INTERNAL} spans with no signal that the agent was stuck
repeating the same search strategy.

The REASONING span kind is essential for characterizing the dominant
failure mode: without it, the 84~reasoning-loop failures appear as
generic ``agent timed out'' events with no diagnostic information
about \emph{why} the agent failed to converge.

\textbf{Closed-loop improvement.}
When the reasoning-loop detector identifies $\geq$3 consecutive
identical calls (via REASONING spans), it injects a strategy-change
prompt.  Table~\ref{tab:closed-loop} shows results across 3~seeds.
At matched iterations, the intervention recovers 9/24 failed instances
(38\%) versus 6/24 for the control (25\%), yielding $+8.3$~pp
attributable solely to the REASONING-based detector.  This closes the
loop from \emph{observability} through \emph{diagnosis} to
\emph{remediation}.

\begin{table}[t]
\centering
\small
\caption{Closed-loop improvement on SWE-bench Lite.  Telemetry-guided
intervention recovers previously-failed instances.}
\label{tab:closed-loop}
\begin{tabular}{@{}lccc@{}}
\toprule
& \textbf{Baseline} & \textbf{Control} & \textbf{Intervention} \\
& \textbf{(8 iter)} & \textbf{(10 iter, no intv)} & \textbf{(10 iter + intv)} \\
\midrule
Recovery (of 24)   & 0/24  & 10/24 (42\%) & 12.7$\pm$1.2/24 (53\%) \\
Combined rate      & 33\%  & 61\%        & 69$\pm$3\% \\
\bottomrule
\end{tabular}
\end{table}

\textbf{Diagnostic quality.}
With \textsc{AgentTelemetry}, developers examine 9.5~spans on average
to reach a full diagnosis versus 28.5~spans with vanilla OTel---a
\textbf{66.8\% reduction} in spans-to-diagnosis.  Fault localization
requires 3~diagnostic spans versus 9.5 (68.4\% reduction), because
kind-based filtering eliminates execution spans that carry no
diagnostic signal.

\textbf{Simulated user study.}
Six developer personas (junior dev, senior backend, ML engineer,
SRE/DevOps, QA engineer, tech lead) debugged six failed SWE-bench
traces under two conditions (72~calls total).  Table~\ref{tab:user-study}
shows that with \textsc{AgentTelemetry}, simulated personas correctly
identified the root cause in 91.7\% of cases (33/36; 95\% CI
[78.2\%, 97.1\%]) versus 25.0\% with vanilla OTel (9/36; 95\% CI
[13.8\%, 41.1\%])---a \textbf{+66.7~pp improvement} (Cohen's $h=1.51$,
large effect).  Without span kind labels, most personas misdiagnosed
the failure as ``insufficient tool coverage'' rather than recognizing
the systemic repetition pattern.  We emphasize this is a
\emph{simulated} study; a controlled human study remains important
future work.

\begin{table}[t]
\centering
\small
\caption{Simulated user study: diagnostic accuracy by persona
($n{=}6$ traces per cell).  AT\,=\,\textsc{AgentTelemetry}.}
\label{tab:user-study}
\begin{tabular}{@{}lcccc@{}}
\toprule
\textbf{Persona} & \textbf{AT Acc.} & \textbf{OTel Acc.} & \textbf{$\Delta$} & \textbf{Exp.} \\
\midrule
Junior Dev        & 5/6 (83\%)  & 1/6 (17\%) & +67\,pp & 2\,yr \\
Sr.\ Backend      & 6/6 (100\%) & 1/6 (17\%) & +83\,pp & 8\,yr \\
ML Engineer       & 5/6 (83\%)  & 1/6 (17\%) & +67\,pp & 5\,yr \\
SRE/DevOps        & 5/6 (83\%)  & 4/6 (67\%) & +17\,pp & 6\,yr \\
QA Engineer       & 6/6 (100\%) & 1/6 (17\%) & +83\,pp & 4\,yr \\
Tech Lead         & 6/6 (100\%) & 1/6 (17\%) & +83\,pp & 10\,yr \\
\midrule
\textbf{Overall}  & \textbf{33/36 (92\%)} & \textbf{9/36 (25\%)} & \textbf{+67\,pp} & --- \\
\bottomrule
\end{tabular}
\end{table}

\subsection{Threats to Validity}

\textbf{Internal.}  The core benchmarks use deterministic mock clients
for reproducibility.  RQ5 complements this with real LLM API validation
across 8~models and 2~providers.  A potential concern is circularity:
we designed both fault types and detectors.  We address this via (1)~85.7\%
coverage of the independently-derived MAST taxonomy~\citep{mast},
(2)~GitHub issue mining confirming all~14 fault types correspond to
real user-reported failures, and (3)~the ablation's structural
argument that is independent of fault design.

\textbf{External.}  Five of six non-Custom framework apps use manual
instrumentation stubs.  To validate real framework code paths, we run
end-to-end tests with LangChain's \texttt{create\_react\_agent},
OpenAI SDK monkey-patching, and a 3-agent multi-agent system---all
producing correct spans and zero false anomalies.

\textbf{Construct.}  FDR is binary; the diagnostic quality analysis
(spans-to-diagnosis: 9.5 vs.\ 28.5) addresses the debugging dimension.

%% ====================================================================
%% 6. RELATED WORK
%% ====================================================================
\section{Related Work}
\label{sec:related}

\textbf{Agent failure taxonomies.}
MAST~\citep{mast} identifies 14~failure modes from 1,600+ annotated
multi-agent traces; 12/14 are detectable via our span kinds.
AgentDebug~\citep{agentdebug} categorizes failures across four
cognitive dimensions.  Aegis~\citep{aegis} classifies agent-environment
failure modes.  AgentOps~\citep{agentops} proposes an observability
taxonomy without an OTel-native implementation or empirical evaluation.
These works characterize \emph{what} goes wrong but do not provide
runtime telemetry to \emph{detect} faults.

\textbf{Debugging tools.}
AgentRx~\citep{agentrx} uses LLM-based analysis of execution logs for
automated multi-agent diagnosis.
AGDebugger~\citep{agdebugger} provides interactive visual debugging,
validated through a 14-person user study.
AgentDiagnose~\citep{agentdiagnose} offers systematic trajectory
analysis.  These tools operate on post-hoc logs;
\textsc{AgentTelemetry} provides the structured telemetry layer they
could consume.

\textbf{Safety and guardrails.}
GuardAgent~\citep{guardagent} and NeMo Guardrails~\citep{nemo_guardrails}
are \emph{enforcement} mechanisms that intercept and block unsafe agent
actions.  Our \texttt{GUARD\_RAIL} span kind serves a complementary
role: it provides \emph{observability} for safety checks, recording
whether a guardrail fired, was bypassed, or failed.  Enforcement
without monitoring creates a blind spot.

\textbf{LLM observability platforms.}
LangSmith~\citep{langsmith} is LangChain-specific.
Langfuse~\citep{langfuse} and Datadog LLM~\citep{datadog_llm} capture
LLM call traces but do not define reusable semantic span kinds.
OpenLLMetry~\citep{openllmetry} and Arize Phoenix intercept API calls
but cannot distinguish agent-level semantics---a planning step and a
summarization call are both \texttt{POST /v1/chat/completions}.

\textbf{OTel standards.}
The GenAI conventions~\citep{otel} define eight operation types;
\textsc{AgentTelemetry} is a strict superset for agent orchestration,
adding five novel span kinds proven necessary by ablation.

%% ====================================================================
%% 7. CONCLUSION
%% ====================================================================
\section{Conclusion}
\label{sec:conclusion}

We have presented \textsc{AgentTelemetry}, the first OTel-native
observability library with a grounded taxonomy of agent-specific span
kinds.  Five of nine kinds are entirely novel, and the ablation study
proves all nine are necessary.  The SWE-bench case study demonstrates
that the novel REASONING span kind characterizes the dominant failure
mode (75\% of instances) and enables a telemetry-guided intervention
improving the patch rate by $+11$~pp.  Controlled fault injection
confirms structural completeness (FDR\,=\,1.000 vs.\ 0.429 for vanilla
OTel).  Instrumentation overhead is negligible ($<$30\,\textmu s per
span).

\textbf{Limitations.}  Core benchmarks use mock LLM clients (validated
with real APIs in RQ5).  Six of seven framework apps use manual
instrumentation stubs.  The SWE-bench study uses a single model and
seed.  The simulated user study uses LLM-generated personas, not human
participants.

Future work includes proposing the taxonomy as an OpenTelemetry
Enhancement Proposal (OTEP), multi-model SWE-bench replication,
and controlled human debugging studies.

%% ====================================================================
%% REFERENCES
%% ====================================================================

\bibliographystyle{plainnat}

\begin{thebibliography}{25}

\bibitem[Cemri et~al.(2025a)]{mast}
M.~Cemri, M.~Z.~Pan, S.~Yang, L.~A.~Agrawal, B.~Chopra, R.~Tiwari,
  K.~Keutzer, A.~Parameswaran, D.~Klein, K.~Ramchandran, M.~Zaharia,
  J.~E.~Gonzalez, and I.~Stoica.
\newblock Why Do Multi-Agent {LLM} Systems Fail?
\newblock In {\em Proc.\ NeurIPS}, 2025.

\bibitem[Cemri et~al.(2025b)]{agentdebug}
M.~Cemri, M.~Y.~Qiang, Y.~Xiao, and L.~Tan.
\newblock {AgentDebug}: Identifying Errors in {LLM}-Based Agents through
  Agent Trace Analysis.
\newblock {\em arXiv preprint arXiv:2504.16744}, 2025.

\bibitem[Chen et~al.(2025)]{aegis}
Z.~Chen, L.~Xu, M.~Wang, et~al.
\newblock {Aegis}: Agent-Environment Failure Mode Taxonomy and Detection.
\newblock {\em arXiv preprint arXiv:2508.19504}, 2025.

\bibitem[{Datadog Inc.}(2024)]{datadog_llm}
{Datadog Inc.}
\newblock {LLM Observability}: Monitor and Improve {LLM} Applications.
\newblock \url{https://www.datadoghq.com/product/llm-observability/}, 2024.

\bibitem[Dong et~al.(2024)]{agentops}
J.~Dong, Y.~Liu, H.~Qian, L.~Zheng, and L.~Chen.
\newblock {AgentOps}: An Observability Taxonomy for {AI} Agent Operations.
\newblock {\em arXiv preprint arXiv:2411.05285}, 2024.

\bibitem[Jiang et~al.(2025)]{agdebugger}
S.~Jiang, X.~Wang, Y.~Zhang, et~al.
\newblock {AGDebugger}: An Interactive Multi-Agent Debugging Tool.
\newblock In {\em Proc.\ ACM CHI}, 2025.

\bibitem[Jimenez et~al.(2024)]{swebench}
C.~E. Jimenez, J.~Yang, A.~Wettig, S.~Yao, K.~Pei, O.~Press, and
  K.~Narasimhan.
\newblock {SWE-bench}: Can Language Models Resolve Real-World {GitHub} Issues?
\newblock In {\em Proc.\ ICLR}, 2024.

\bibitem[{LangChain Inc.}(2024)]{langsmith}
{LangChain Inc.}
\newblock {LangSmith}: Platform for Building Production-Grade {LLM}
  Applications.
\newblock \url{https://smith.langchain.com/}, 2024.

\bibitem[{Langfuse GmbH}(2024)]{langfuse}
{Langfuse GmbH}.
\newblock {Langfuse}: Open-Source {LLM} Observability Platform.
\newblock \url{https://langfuse.com/}, 2024.

\bibitem[Li et~al.(2025)]{agentdiagnose}
Y.~Li, Z.~Wang, H.~Chen, et~al.
\newblock {AgentDiagnose}: A Toolkit for Diagnosing Agent Trajectories.
\newblock In {\em Proc.\ EMNLP (Demo Track)}, 2025.

\bibitem[Liu et~al.(2024)]{agentbench}
X.~Liu, H.~Yu, H.~Zhang, et~al.
\newblock {AgentBench}: Evaluating {LLMs} as Agents.
\newblock In {\em Proc.\ ICLR}, 2024.

\bibitem[Ma et~al.(2024)]{agentboard}
C.~Ma, J.~Zhang, Z.~Zhu, C.~Yang, Y.~Yang, et~al.
\newblock {AgentBoard}: An Analytical Evaluation Board of Multi-Turn {LLM}
  Agents.
\newblock In {\em Proc.\ ACL}, 2024.

\bibitem[{Microsoft Research}(2026)]{agentrx}
{Microsoft Research}.
\newblock {AgentRx}: Automated Diagnostic Framework for Multi-Agent Systems.
\newblock {\em arXiv preprint arXiv:2602.02475}, 2026.

\bibitem[{OpenTelemetry Authors}(2024)]{otel}
{OpenTelemetry Authors}.
\newblock {OpenTelemetry Specification}, v1.30.
\newblock \url{https://opentelemetry.io/docs/specs/otel/}, 2024.

\bibitem[Strauss and Corbin(1998)]{strauss_corbin}
A.~Strauss and J.~Corbin.
\newblock {\em Basics of Qualitative Research: Techniques and Procedures for
  Developing Grounded Theory}.
\newblock Sage, 2nd edition, 1998.

\bibitem[{Traceloop}(2024)]{openllmetry}
{Traceloop}.
\newblock {OpenLLMetry}: Open-Source Observability for {LLM} Applications
  Built on {OpenTelemetry}.
\newblock \url{https://github.com/traceloop/openllmetry}, 2024.

\bibitem[Zhan et~al.(2025)]{guardagent}
X.~Zhan, Y.~Luo, S.~Wang, et~al.
\newblock {GuardAgent}: Safeguard {LLM} Agents by a Guard Agent via
  Knowledge-Enabled Reasoning.
\newblock In {\em Proc.\ ICML}, 2025.

\bibitem[Rebedea et~al.(2023)]{nemo_guardrails}
T.~Rebedea, R.~Dinu, M.~Sreedhar, C.~Parisien, and J.~Cohen.
\newblock {NeMo Guardrails}: A Toolkit for Controllable and Safe {LLM}
  Applications.
\newblock In {\em Proc.\ EMNLP (Demo Track)}, 2023.

\end{thebibliography}

%%%%%%%%%%%%%%%%%%%%%%%%%%%%%%%%%%%%%%%%%%%%%%%%%%%%%%%%%%%%

\appendix

\section{Additional Results}

\textbf{MAST mapping.}
Table~\ref{tab:mast-mapping} maps MAST's 14~failure modes to our fault types.  12 of~14 (85.7\%) are detectable; the two gaps (FM-2.2: failure to ask for clarification, FM-3.2: missing verification) are omission failures requiring expected-span specifications.

\begin{table}[ht]
\centering
\scriptsize
\caption{Mapping MAST failure modes to AgentTelemetry fault types.}
\label{tab:mast-mapping}
\begin{tabular}{@{}llll@{}}
\toprule
\textbf{MAST Mode} & \textbf{Code} & \textbf{Our Fault Type(s)} & \textbf{Cov.} \\
\midrule
Disobey Task Spec   & FM-1.1 & wrong\_tool, hallucination  & Full \\
Disobey Role Spec   & FM-1.2 & agent\_misroute, GR bypass  & Full \\
Step Repetition     & FM-1.3 & infinite\_loop, reas.\ loop & Full \\
Loss of History     & FM-1.4 & ctx.\ overflow, mem.\ corr. & Full \\
Unaware of Term.    & FM-1.5 & infinite\_loop, plan.\ fail & Full \\
Conversation Reset  & FM-2.1 & mem.\ corr., ctx.\ overflow & Partial \\
Fail to Clarify     & FM-2.2 & --- & Gap \\
Task Derailment     & FM-2.3 & wrong\_tool, plan.\ fail    & Partial \\
Info Withholding    & FM-2.4 & circular\_delegation         & Partial \\
Ignored Input       & FM-2.5 & agent\_misroute, circ.\ del. & Partial \\
Reas.-Action Mism.  & FM-2.6 & wrong\_tool, hallucination  & Full \\
Premature Term.     & FM-3.1 & timeout, plan.\ fail        & Partial \\
No Verification     & FM-3.2 & --- & Gap \\
Incorrect Verif.    & FM-3.3 & GR bypass, hallucination    & Full \\
\bottomrule
\end{tabular}
\end{table}

%%%%%%%%%%%%%%%%%%%%%%%%%%%%%%%%%%%%%%%%%%%%%%%%%%%%%%%%%%%%

\newpage
\section*{NeurIPS Paper Checklist}

\begin{enumerate}

\item {\bf Claims}
    \item[] Question: Do the main claims made in the abstract and introduction accurately reflect the paper's contributions and scope?
    \item[] Answer: \answerYes{}
    \item[] Justification: The abstract and introduction state six specific contributions (taxonomy, controlled benchmark, ablation, head-to-head comparison with OpenLLMetry, SWE-bench case study with statistically-disclosed intervention result, and open-source library), each supported by a corresponding experimental research question (RQ1--RQ6) in \S5.

\item {\bf Limitations}
    \item[] Question: Does the paper discuss the limitations of the work performed by the authors?
    \item[] Answer: \answerYes{}
    \item[] Justification: \S7 (Conclusion) includes an explicit Limitations paragraph with four substantive limitations: (L1) mock-vs-real evaluation gap, (L2) closed-loop intervention sample size and Fisher's exact $p{=}0.53$ at $n{=}24$, (L3) LLM-judged plausibility vs.\ official SWE-bench harness verification, (L4) simulated user study label-visibility confound.

\item {\bf Theory assumptions and proofs}
    \item[] Question: For each theoretical result, does the paper provide the full set of assumptions and a complete (and correct) proof?
    \item[] Answer: \answerNA{}
    \item[] Justification: The paper does not include theoretical results or proofs. All claims are empirically validated.

    \item {\bf Experimental result reproducibility}
    \item[] Question: Does the paper fully disclose all the information needed to reproduce the main experimental results of the paper to the extent that it affects the main claims and/or conclusions of the paper (regardless of whether the code and data are provided or not)?
    \item[] Answer: \answerYes{}
    \item[] Justification: \S5 specifies all experimental parameters (14 fault types, 5 conditions, 7 frameworks, 13 real-LLM models, 2{,}940 cross-product cells reducing to 490 independent-signal cells under deterministic mocks). The SWE-bench setup details model, iteration count, tool set, and cost. All controlled benchmarks use deterministic mock clients for bit-exact reproducibility. Random seeds fixed at $\{42, 123, 456, 789, 1024\}$ for multi-seed analysis.

\item {\bf Open access to data and code}
    \item[] Question: Does the paper provide open access to the data and code, with sufficient instructions to faithfully reproduce the main experimental results, as described in supplemental material?
    \item[] Answer: \answerYes{}
    \item[] Justification: The library is published on PyPI (\texttt{pip install agenttelemetry}) under Apache-2.0 license. Source code at the public GitHub repository (linked in supplementary). Dataset hosted on Hugging Face Datasets with Croissant 1.0 metadata file (\texttt{metadata/croissant.json}) including all mandatory Responsible AI fields. License: Apache-2.0 for code, CC-BY-4.0 for traces.

\item {\bf Experimental setting/details}
    \item[] Question: Does the paper specify all the training and test details (e.g., data splits, hyperparameters, how they were chosen, type of optimizer) necessary to understand the results?
    \item[] Answer: \answerYes{}
    \item[] Justification: \S5 provides complete experimental setup including models, frameworks, fault injection rates, privacy levels, iteration counts, SWE-bench instance selection criteria, and detector thresholds. \S5.7 reports a full $\pm{}50\%$ threshold sensitivity sweep with provenance for each default value.

\item {\bf Experiment statistical significance}
    \item[] Question: Does the paper report error bars suitably and correctly defined or other appropriate information about the statistical significance of the experiments?
    \item[] Answer: \answerYes{}
    \item[] Justification: RQ1 reports FDR with standard deviations across 5 random seeds at 5 injection rates. The closed-loop intervention reports an exact Fisher's two-sided p-value ($p{=}0.53$ at $n{=}24$) and explicitly discloses non-significance plus the power calculation ($\sim{}214$ instances per arm needed at observed effect). FPR$=$0.000 across all 10 false-positive runs and separately confirmed on 112 clean SWE-bench traces. The simulated user study reports Cohen's $h$ with explicit framing as a label-visibility upper bound.

\item {\bf Experiments compute resources}
    \item[] Question: For each experiment, does the paper provide sufficient information on the computer resources (type of compute workers, memory, time of execution) needed to reproduce the experiments?
    \item[] Answer: \answerYes{}
    \item[] Justification: Overhead benchmarks specify Apple M4 Pro hardware. SWE-bench cost reported (\$2.53). Real-LLM corpus cost reported (\$0.62 across 13 models). Larger-sample replication script ($n{=}60$ via Opus 4.6) released with cost estimate \$10--20.

\item {\bf Code of ethics}
    \item[] Question: Does the research conducted in the paper conform, in every respect, with the NeurIPS Code of Ethics?
    \item[] Answer: \answerYes{}
    \item[] Justification: The research involves no human subjects, no personally identifiable data, and no dual-use concerns. The library is designed to improve system reliability and transparency. Privacy controls explicitly mitigate the only plausible negative externality (chain-of-thought capture).

\item {\bf Broader impacts}
    \item[] Question: Does the paper discuss both potential positive societal impacts and negative societal impacts of the work performed?
    \item[] Answer: \answerYes{}
    \item[] Justification: An explicit Broader Impact section discusses positive externalities (faster fault diagnosis, runtime guardrails, common vocabulary) and the only material negative externality (REASONING/PLANNING spans capture chain-of-thought potentially containing user PII), with the three-level privacy model as concrete mitigation.

\item {\bf Safeguards}
    \item[] Question: Does the paper describe safeguards that have been put in place for responsible release of data or models that have a high risk for misuse?
    \item[] Answer: \answerYes{}
    \item[] Justification: The released traces use synthetic fault-injection inputs and contain no human-subject data. The default \texttt{METADATA\_ONLY} privacy level captures no prompt content (sufficient for fault detection: FDR is identical at metadata and full capture levels).

\item {\bf Licenses for existing assets}
    \item[] Question: Are the creators or original owners of assets (e.g., code, data, models), used in the paper, properly credited and are the license and terms of use explicitly mentioned and properly respected?
    \item[] Answer: \answerYes{}
    \item[] Justification: SWE-bench (\texttt{princeton-nlp/SWE-bench\_Lite}, MIT-licensed) is properly cited. OpenTelemetry (Apache-2.0) is cited. All seven framework adapters (LangChain, CrewAI, AutoGen, LlamaIndex, Anthropic SDK, OpenAI SDK, Custom) are referenced with their respective licenses. MAST and AgentRx datasets/papers cited.

\item {\bf New assets}
    \item[] Question: Are new assets introduced in the paper well documented and is the documentation provided alongside the assets?
    \item[] Answer: \answerYes{}
    \item[] Justification: The AgentTelemetry library is documented with API reference, installation instructions, and usage examples in the repository. The span taxonomy is fully specified in the paper. The dataset includes a Datasheet for Datasets (\S Appendix) and Croissant 1.0 metadata.

\item {\bf Crowdsourcing and research with human subjects}
    \item[] Question: For crowdsourcing experiments and research with human subjects, does the paper include the full text of instructions given to participants and screenshots, if applicable, as well as details about compensation (if any)?
    \item[] Answer: \answerNA{}
    \item[] Justification: The paper does not involve crowdsourcing or human subjects. The simulated user study uses LLM-generated personas (GPT-4o-mini role-play), explicitly disclosed and explicitly framed as a label-visibility upper bound rather than a substitute for a controlled human study. A controlled human study is queued for the journal extension.

\item {\bf Institutional review board (IRB) approvals or equivalent for research with human subjects}
    \item[] Question: Does the paper describe potential risks incurred by study participants, whether such risks were disclosed to the subjects, and whether Institutional Review Board (IRB) approvals (or an equivalent approval/review based on the requirements of your country or institution) were obtained?
    \item[] Answer: \answerNA{}
    \item[] Justification: No human subjects were involved. The simulated user study uses LLM-generated role-plays, not human participants.

\item {\bf Declaration of LLM usage}
    \item[] Question: Does the paper describe the usage of LLMs if it is an important, original, or non-standard component of the core methods in this research?
    \item[] Answer: \answerYes{}
    \item[] Justification: LLMs are used as experimental subjects (GPT-4o-mini for SWE-bench, 13 models from Anthropic and OpenAI for real-API validation, Opus 4.6 for queued $n{=}60$ matched-control extension) and for the simulated user study (GPT-4o-mini for persona role-plays, framed as label-visibility upper bound). All LLM uses are described in \S5.5--5.6 and \S Limitations.

\end{enumerate}



\end{document}
